\subsection{结论：输入效率改善，稳定提分仍需验证}
S两次共同评分为0.5846与0.6077，均值0.5962，恰好等于前次两条balanced
固定答案在本批评委下的均值。同期新B为0.5308，显示对照自身也有明显运行波动。
因此不能只选S较高的一次或只比较较低的同期B，就宣称稳定全面提分。
前次原评分仍保留；本节使用的是新共同评委重评，不是改写旧答案。

E为0.5846，65题均有有效回答，但尚无独立重复；其重复因额度16\%低于
预设启动门槛而未执行，并非因分数低而取消。
E较同期B的QA总输入减少38.2\%，未缓存输入减少63.6\%；
较S两次平均分别再减少15.1\%和27.5\%。
更稳妥的定位是\textbf{保持此前质量附近的低输入候选}，而不是已经证实的无损压缩或
稳定准确率提升。建议保留balanced作回归基线，在新额度窗口先完成E重复。

\subsection{三条改进路线的实际边界}
仅原文事实窗口F为0.5000，说明切短本身不保证效果。
S增加了来源事实向量索引等组合变化；其371次实际查询离线重放全部复现原入选集合，
355次在关闭索引后主片段选择改变，19次选入原候选池之外的节点。
因此多数变化是重排与细粒度准入，不应写成大幅扩大正确证据召回。
重放未使用gold，只说明检索路径，不证明提分因果。

FI相对F总分差为$-0.0077$，探索性区间跨零。
其人物理解类别有所改善，但知识提取下降；类别重叠，不能挑有利子集代表整体。
当前name/role匹配仍是来源声明，不是人脸、声音和姓名的核实绑定。
部分错例已经看到了正确词，却把偏好或物品归给了错误对象；
解决这种错误需要可信的实体--事实关系，而非更强的名字共现过滤。

首版增量计数会把不发送给模型的审计字段算入基数：
1000代理单位反例可准入7095，修复后有无该字段均只增加890。
\textbf{修复前总输入上限仍有效}，这是新增额度失真，而非旧实验输入越界。
E改为实际模型投影计数和3500步长，370次QA调用没有epoch重建；
305次相邻消息前缀全部保持，实际最大新增量3500。
消息前缀检查不包含工具schema，也不等价于供应商缓存保证。

57次请求设置required且仅提供Answer，其中41次返回原生Answer，16次仍返回普通文本，
由既有有效文本兼容路径接受，没有追加调用。
不能把该参数描述成严格原生输出保证。45秒回答预留本轮也未触发。
此外，32节点/每节点12条约束的是本次新增准入，不是整个增量历史；
E最多可见38个节点，全局仍以总输入预算约束。

\subsection{与M3-Agent的取舍}
M3参考适配两次轨迹在本批评委下均为0.4615，逐题答案并不相同。
S均值高0.1346，E高0.1231，但E的QA输入仍为M3两次平均的5.20倍，
未缓存输入约7.77倍。M3的紧凑事件/语义事实、身份条件检索与少量新片段交付
仍有明显效率优势，不能从本批得分推出我们的身份图推理更强。
两者控制协议和末轮猜测策略不同，且参考方是同模型适配，不是官方checkpoint复现。

\subsection{审计结果与用量采集缺口}
七个新oGMemory轨迹共455题，454个有效回答，唯一失败为FIP的Search文本协议；
M3同期重复65题均有效。251,023个实际可见原文位置经还原验证，
数据库内容与事实索引hash不变。来源一致不保证源记忆和最终答案真实正确。
本轮未重新生成AV感知，仍受前缀帧采样、ASR与人物归属问题影响；
current\_clip允许看至当前片段末尾，非严格因果在线评测。

\textbf{共同评委用量日志存在缺口：}
715条评分已完成且无评委错误，但Python启动时sitecustomize因源码路径缺失，
未成功安装用量采集。714个需要模型评分的任务没有持久化usage或逐响应模型标识；
已结束进程的内存计数无法恢复。配置与冻结代码仍可核查。
没有重评成功项或用估计值补账。全部QA轨迹用量完整，故QA输入比较仍有效，
但不能给出整轮实验的完整总token或费用。
已修复后续启动环境并增加采集未就绪即禁止模型调用的回归检查。

工程验证包括1447项扩展检查通过、108项跳过及9项此前已复现的既有失败排除；
根仓评测与审计工具79项测试通过。排除并不意味着旧失败已修复。
新增索引构建入口使用只读事务和租户/视频过滤，真实数据dry-run及原向量写入回放通过，
未覆盖本轮固定索引。

本研究仅5个开发视频、65题、同模型自动评委。聚类重采样区间只刻画固定答案上的
视频波动，不能覆盖所有生成与评委随机性，亦非多重开发比较后的确认性检验。
详细证据与未运行组状态见
\code{docs/m3_memory_design_findings.md}和
\code{experiments/m3_memory_design_20261102/}。
